\documentclass[11pt]{amsart}

\usepackage{amsmath,amssymb,amsthm,mathtools}
\usepackage{microtype}
\usepackage[hidelinks]{hyperref}
\usepackage{url}


\renewcommand{\subjclassname}{2020 Mathematics Subject Classification}

\numberwithin{equation}{section}
\setlength{\emergencystretch}{3em}
\urlstyle{same}
\raggedbottom

\DeclareMathOperator{\Log}{Log}
\DeclareMathOperator{\Arg}{Arg}

\newtheorem{theorem}{Theorem}[section]
\newtheorem{proposition}[theorem]{Proposition}
\newtheorem{lemma}[theorem]{Lemma}
\newtheorem{corollary}[theorem]{Corollary}
\theoremstyle{definition}
\newtheorem{definition}[theorem]{Definition}
\theoremstyle{remark}
\newtheorem{remark}[theorem]{Remark}

\newtheorem*{theoremA}{Theorem A}
\newtheorem*{corollaryB}{Corollary B}
\newtheorem*{corollaryC}{Corollary C}

\newcommand{\C}{\mathbb C}
\newcommand{\R}{\mathbb R}
\newcommand{\D}{\mathcal D}
\newcommand{\Hplus}{\mathbb H_+}

\title[Radial uniqueness of the nonzero 1-points of $J_0$]
{Radial Uniqueness of the Nonzero \texorpdfstring{$1$}{1}-Points of \texorpdfstring{$J_0$}{J0}}
\author{Helen W.~J. Zhang}
\address{School of Mathematics, Hunan University}
\email{helenzhang@hnu.edu.cn}
\thanks{Helen W.~J. Zhang was supported by the National Natural Science Foundation of China (Grant No.~12371327) and Hunan Provincial Natural Science Foundation of China (Grant No.~2026JJ50107).}

\author{Chengdong Zhao}
\address{School of Mathematics and Statistics, Central South University}
\email{cdzhao@csu.edu.cn}
\date{}
\subjclass[2020]{Primary 33C10; Secondary 30D20, 31A05, 35J05}
\keywords{Bessel function, 1-point, radial uniqueness, radial escape, circular mean, two-radius theorem}

\begin{document}

\begin{abstract}
We answer Problem~2.45 in the Hayman--Lingham problem collection.  If
$z_0\ne0$, $J_0(z_0)=1$, and $c>1$, then
$|J_0(cz_0)|>c^{7/25}$.  Hence no two distinct nonzero
$1$-points of $J_0$ lie on the same open ray from the origin.  The proof uses
an exact, branch-consistent Laplace decomposition into growing and decaying
components, and uniform estimates show that
the separation persists along the outward ray and yields a positive radial
logarithmic derivative.  As an application, for distinct positive radii
$r_1,r_2$ and $R>r_1+r_2$, raywise uniqueness makes Volchkov's positive nonunit
spectral-ratio exclusion automatic in the local planar spherical-mean
theorem and gives the corresponding global $L^1_{\mathrm{loc}}$ result.
\end{abstract}

\maketitle

\section{Introduction}

Let $J_0$ denote the Bessel function of the first kind of order zero.
Problem~2.45 in the Hayman--Lingham edition of Hayman's problem collection
asks whether $J_0(z)=1$ has at most one nonzero solution on each ray from
the origin.  Hayman and Lingham attribute the question to L.~Zalcman and
record that the available asymptotic information gives only finitely many
solutions on any fixed ray
\cite[Problem~2.45]{HaymanLingham2019}.

We answer this question by proving the following stronger estimate.

\begin{theoremA}
If $z_0\in\C\setminus\{0\}$ and $J_0(z_0)=1$, then
\[
   |J_0(cz_0)|>c^{7/25}\qquad(c>1).
\]
\end{theoremA}

Since $c^{7/25}>1$ for $c>1$, Theorem~A answers Problem~2.45 immediately.

\begin{corollaryB}
For every $\theta\in[0,2\pi)$,
\[
   \#\{r>0:J_0(re^{i\theta})=1\}\le1.
\]
Equivalently, no two distinct nonzero $1$-points of $J_0$ lie on the same
open ray from the origin.
\end{corollaryB}

The question also has a spectral interpretation in converse mean-value
theorems.  Delsarte and Lions showed that two fixed spherical mean
identities force harmonicity when the corresponding averaging operators
have no nonzero common fixed spectral mode \cite{DelsarteLions1959}; the
Euclidean formulation is given in Delsarte's lectures
\cite[Part~III, \S\S3--4]{Delsarte1961}.  In dimension two the common
spectral equations are
\[
   J_0(i\lambda r_1)=J_0(i\lambda r_2)=1,\qquad\lambda\ne0.
\]
If $0<r_1<r_2$, $c=r_2/r_1$, and $z=i\lambda r_1$, these equations become
$J_0(z)=J_0(cz)=1$.  Thus a positive nonunit spectral ratio occurs exactly
when two nonzero $1$-points lie on one open ray.

Volchkov's local planar classification separates the strict geometric
condition $R>r_1+r_2$ from the exclusion of a Bessel root-ratio set
\cite{Volchkov1996}.  Theorem~A removes the positive nonunit part of this
spectral obstruction, yielding the following consequence.

\begin{corollaryC}
Let $B_R:=\{x\in\R^2:|x|<R\}$, let $r_1,r_2$ be distinct positive radii,
and let $R>r_1+r_2$.  If $f\in L^1_{\mathrm{loc}}(B_R)$ satisfies the
normalized circular mean identities at $r_1$ and $r_2$ almost everywhere
on their natural domains, then $\Delta f=0$ in $\D'(B_R)$.  In the
full-space case, if $f\in L^1_{\mathrm{loc}}(\R^2)$ satisfies the two
identities almost everywhere on $\R^2$, then $\Delta f=0$ in $\D'(\R^2)$.
\end{corollaryC}

A radial energy identity and an elementary zero-free disk first localize
every nonzero $1$-point.  On the
principal branch of $K_0$, a Laplace representation then gives an exact,
branch-consistent decomposition of $J_0$ into growing and decaying
components.  Uniform amplitude estimates separate the two components at
each principal-sector $1$-point, and this separation persists along the
whole outward ray.  The resulting lower bound for the radial logarithmic
derivative implies Theorem~A; evenness and conjugation symmetry give
the global statement.

This paper was submitted on 30 August 2026.  Ferudun's preprint
\cite{Ferudun2026}, dated 3 October 2026, presents a computer-assisted proof
of the same raywise uniqueness statement.  Both arguments separate growing
and decaying Bessel components.  Ferudun combines explicit Hankel bounds,
contraction and Rouch\'e arguments, certified root counting, and ordering
of root arguments.  Our proof uses an exact $K_0$ Laplace decomposition and
uniform amplitude and derivative estimates propagated along outward rays
to obtain the quantitative estimate of Theorem~A, without enumerating the
$1$-points.

Section~\ref{sec:localization} gives the localization.  Sections
\ref{sec:laplace} and~\ref{sec:components} establish the exact decomposition
and the uniform component bounds.  Section~\ref{sec:escape} proves radial
escape and answers Problem~2.45.  Section~\ref{sec:two-radius} proves the
two-radius consequences and compares them with Volchkov's local theorem.

The results proved here have been formalized in Lean~4 using
mathlib~\cite{Lean4,mathlib}. The source code, statement correspondence,
and reproducible build instructions are available in the
\href{https://github.com/VerMoriarty/j0-radial-uniqueness}
{accompanying GitHub repository}.

\section{Preliminary localization of the nonzero \texorpdfstring{$1$}{1}-points}
\label{sec:localization}

We use the defining series
\begin{equation*}
   J_0(z):=\sum_{m=0}^{\infty}\frac{(-1)^m z^{2m}}{4^m(m!)^2},
   \qquad z\in\C.
\end{equation*}
Set
\[
   F(w):=\sum_{m=0}^{\infty}\frac{w^m}{(m!)^2},
   \qquad
   Q(w):=\sum_{n=0}^{\infty}\frac{w^n}{((n+1)!)^2}.
\]
Then
\begin{equation}\label{eq:J0-FQ-relations}
   J_0(z)=F\!\left(-\frac{z^2}{4}\right),
   \qquad
   F(w)-1=wQ(w).
\end{equation}

\begin{lemma}\label{lem:F-right-half-plane}
If $w\ne0$ and $\Re w\ge0$, then $|F(w)|>1$.
\end{lemma}

\begin{proof}
Termwise differentiation gives $wF''(w)+F'(w)=F(w)$.  Put
$h(s)=F(sw)$ and $Y(s)=|h(s)|^2$.  Then $sh''+h'=wh$, and integration of
$(sY')'=2(s|h'|^2+\Re w\,|h|^2)$ gives
\begin{equation}\label{eq:F-energy-identity}
   sY'(s)=2\int_0^s
   \left(u|h'(u)|^2+\Re w\,|h(u)|^2\right)du.
\end{equation}
The integrand is nonnegative for $u\ge0$.  Since $h'(0)=w\ne0$, its first
term is positive for all sufficiently small $u>0$.  Thus $Y'(s)>0$ for
every $s>0$, and $|F(w)|^2=Y(1)>Y(0)=1$.
\end{proof}

\begin{lemma}\label{lem:no-real-one-points}
For every $x\in\R\setminus\{0\}$, one has $|J_0(x)|<1$.
\end{lemma}

\begin{proof}
Put $q=x^2/4>0$ and $h(s)=F(-qs)$.  The equation for $F$ gives
$sh''+h'=-qh$.  Hence the energy
\[
   E(s):=|h(s)|^2+\frac{s}{q}|h'(s)|^2
   \qquad\text{satisfies}\qquad
   E'(s)=-\frac1q|h'(s)|^2.
\]
Since $h'(0)=-q\ne0$, integration on $[0,1]$ gives $E(1)<E(0)=1$.
Therefore $|J_0(x)|^2=|h(1)|^2\le E(1)<1$.
\end{proof}

\begin{lemma}\label{lem:Q-zero-free}
The function $Q$ has no zero in the disk $|w|<4$.
\end{lemma}

\begin{proof}
Write $w=4u$, where $|u|<1$, and put $b_n=4^n/((n+1)!)^2$.  The sequence
$(b_n)$ is nonnegative and summable, with $b_0=1$ and
$b_{n+1}/b_n=4/(n+2)^2\le1$.  Set $d_n=b_n-b_{n+1}\ge0$.  Absolute
convergence and telescoping give
\begin{equation}\label{eq:Q-Abel-transform}
   (1-u)Q(4u)=1-\sum_{n=0}^{\infty}d_nu^{n+1},
   \qquad
   \sum_{n=0}^{\infty}d_n=1.
\end{equation}
Moreover,
\[
   \left|\sum_{n=0}^{\infty}d_nu^{n+1}\right|
   \le |u|\sum_{n=0}^{\infty}d_n=|u|<1.
\]
Thus $(1-u)Q(4u)\ne0$, and hence $Q(w)\ne0$.
\end{proof}

\begin{proposition}\label{prop:global-localization}
If $J_0(z)=1$ and $z\ne0$, then
$
   \Re(z^2)>0$, $z\notin\R$, $|z|\ge4$.
\end{proposition}

\begin{proof}
Let $w=-z^2/4$.  If $\Re(z^2)\le0$, then $w\ne0$ and $\Re w\ge0$, so
Lemma~\ref{lem:F-right-half-plane} gives $|J_0(z)|=|F(w)|>1$, a
contradiction.  Lemma~\ref{lem:no-real-one-points} excludes nonzero real
$z$.  Finally, \eqref{eq:J0-FQ-relations} gives $wQ(w)=0$.  Since $w\ne0$,
Lemma~\ref{lem:Q-zero-free} yields $|w|\ge4$, and hence $|z|\ge4$.
\end{proof}

Thus every nonzero $1$-point lies in one of the four open sectors adjacent
to the real axis.  By evenness and conjugation symmetry, it will later be
enough to work in $0<\arg z<\pi/4$.

\section{An exact Laplace decomposition}\label{sec:laplace}

Let $\Omega=\C\setminus(-\infty,0]$.  On $\Omega$ we use the principal
logarithm $\Log w=\log|w|+i\Arg w$, where $-\pi<\Arg w<\pi$, and define all
half-integral powers through this branch.  Let $K_0$ denote the modified
Bessel function of the second kind of order zero, with its principal branch.
Put
\[
   d\mu(t):=\frac{e^{-t}t^{-1/2}}{\sqrt\pi}\,dt,
   \qquad
   \int_0^\infty d\mu(t)=1,
   \qquad
   \int_0^\infty t\,d\mu(t)=\frac12.
\]
For $w\in\Omega$, define
\[
   \mathcal I(w):=\int_0^\infty
   \left(1+\frac{t}{2w}\right)^{-1/2}d\mu(t),
   \qquad
   \mathcal K(w):=\sqrt{\frac{\pi}{2w}}e^{-w}\mathcal I(w).
\]

\begin{proposition}\label{prop:K-Laplace}
The functions $\mathcal I$ and $\mathcal K$ are holomorphic on $\Omega$,
\[
   \mathcal I'(w)=\frac{1}{4w^2}\int_0^\infty
   t\left(1+\frac{t}{2w}\right)^{-3/2}d\mu(t),
\]
and $\mathcal K(w)=K_0(w)$ for every $w\in\Omega$.
\end{proposition}

\begin{proof}
For $w\in\Omega$ and $t>0$, the point $1+t/(2w)$ does not lie on
$(-\infty,0]$: otherwise $1+t/(2w)=-s$ for some $s\ge0$, and then
$w=-t/[2(1+s)]\in(-\infty,0)$.  Thus the integrand is holomorphic in $w$ on
$\Omega$.

Let $L\Subset\Omega$ be nonempty and compact.  Compactness gives positive
constants $m_L,M_L,\eta_L$ such that, for $w\in L$ and $0\le s\le1$,
\[
   m_L\le |w|\le M_L,
   \qquad |2(1-s)w+s|\ge\eta_L.
\]
Indeed, a zero of the latter expression would force either $s=1$ and
$1=0$, or $w=-s/[2(1-s)]\in(-\infty,0]$.  Taking $s=t/(1+t)$ shows that
\[
   \left|1+\frac{t}{2w}\right|\ge c_L:=\frac{\eta_L}{2M_L}>0
   \qquad(w\in L,\ t\ge0).
\]
The kernel is therefore bounded on $L$ by $c_L^{-1/2}$.  Its derivative is
\[
   \frac{\partial}{\partial w}
   \left(1+\frac{t}{2w}\right)^{-1/2}
   =\frac{t}{4w^2}\left(1+\frac{t}{2w}\right)^{-3/2},
\]
whose modulus is at most $t c_L^{-3/2}/(4m_L^2)$.  The mass and first
moment of $\mu$ now justify differentiation under the integral sign.
This proves the holomorphy of $\mathcal I$ and its derivative formula;
the defining prefactor then gives the holomorphy of $\mathcal K$.

For $w>0$, the standard representation \cite[10.32.9]{DLMF}
\[
   K_0(w)=\int_0^\infty e^{-w\cosh u}\,du
\]
and the substitution $t=w(\cosh u-1)$ give
\[
   K_0(w)=e^{-w}\int_0^\infty
   \frac{e^{-t}}{\sqrt{t(2w+t)}}\,dt=\mathcal K(w).
\]
The principal branch of $K_0$ is holomorphic on $\Omega$
\cite[Sec.~10.25(ii)]{DLMF}.  Since $\Omega$ is connected, the identity
theorem yields $\mathcal K=K_0$ throughout $\Omega$.
\end{proof}

Let $\Hplus=\{z\in\C:\Re z>0\}$.  For $z\in\Hplus$, define
\begin{equation*}
   A_\pm(z):=\int_0^\infty
   \left(1\pm\frac{it}{2z}\right)^{-1/2}d\mu(t).
\end{equation*}

\begin{proposition}\label{prop:exact-decomposition}
The functions $A_\pm$ are holomorphic on $\Hplus$ and satisfy
\begin{equation}\label{eq:Apm-derivative}
   A_\pm'(z)=\pm\frac{i}{4z^2}\int_0^\infty
   t\left(1\pm\frac{it}{2z}\right)^{-3/2}d\mu(t).
\end{equation}
Moreover, for every $z\in\Hplus$,
\begin{equation}\label{eq:exact-J0-decomposition}
   J_0(z)=\frac{e^{i(z-\pi/4)}A_+(z)+e^{-i(z-\pi/4)}A_-(z)}
   {\sqrt{2\pi z}},
\end{equation}
where the square root is taken on the principal branch.
\end{proposition}

\begin{proof}
Since $-iz,iz\in\Omega$ for $z\in\Hplus$, one has
$A_+(z)=\mathcal I(-iz)$ and $A_-(z)=\mathcal I(iz)$.  Holomorphy and
\eqref{eq:Apm-derivative} follow from Proposition~\ref{prop:K-Laplace} and
the chain rule.

Write $z=re^{i\theta}$ with $-\pi/2<\theta<\pi/2$.  Then
\[
   \Log(iz)=\Log z+\frac{\pi i}{2},
   \qquad
   \Log(-iz)=\Log z-\frac{\pi i}{2},
\]
and therefore
\begin{equation*}
   (iz)^{-1/2}=e^{-\pi i/4}z^{-1/2},
   \qquad
   (-iz)^{-1/2}=e^{\pi i/4}z^{-1/2}.
\end{equation*}
Proposition~\ref{prop:K-Laplace} gives
\[
\begin{aligned}
   K_0(-iz)&=\sqrt{\frac\pi2}\,e^{\pi i/4}z^{-1/2}e^{iz}A_+(z),\\
   K_0(iz)&=\sqrt{\frac\pi2}\,e^{-\pi i/4}z^{-1/2}e^{-iz}A_-(z).
\end{aligned}
\]
The standard connection formula \cite[10.27.9]{DLMF}, specialized to order
zero, is
\[
   \pi iJ_0(z)=K_0(-iz)-K_0(iz),
   \qquad \Re z>0,
\]
with both values taken on the principal branch.  Substitution and the identities
\[
   \frac{1}{\pi i}\sqrt{\frac\pi2}e^{\pi i/4}
   =\frac{e^{-\pi i/4}}{\sqrt{2\pi}},
   \qquad
   -\frac{1}{\pi i}\sqrt{\frac\pi2}e^{-\pi i/4}
   =\frac{e^{\pi i/4}}{\sqrt{2\pi}}
\]
yield \eqref{eq:exact-J0-decomposition}.
\end{proof}

\section{Control of the two components}\label{sec:components}

Let $z=x+iy=\rho e^{i\theta}$, where $\rho>0$ and
$0\le\theta\le\pi/4$, and put $\kappa=2^{-9/4}$.

\begin{proposition}\label{prop:A-uniform-estimates}
In this closed sector,
\begin{equation}\label{eq:A-uniform-estimates}
   |A_\pm(z)-1|\le\frac{\kappa}{\rho},
   \qquad
   |A_\pm'(z)|\le\frac{\kappa}{\rho^2}.
\end{equation}
If $\rho\ge4$, then
\begin{equation}\label{eq:A-modulus-bounds}
   \frac{15}{16}<|A_\pm(z)|<\frac{17}{16},
   \qquad
   \left|z\frac{A_\pm'(z)}{A_\pm(z)}\right|<\frac1{15}.
\end{equation}
\end{proposition}

\begin{proof}
Since $x\ge y\ge0$ and $x>0$, completing squares in Cartesian
coordinates gives, for $t\ge0$ and $0\le\sigma\le1$,
\[
   \left|1+\sigma\frac{it}{2z}\right|^2\ge1,
   \qquad
   \left|1-\sigma\frac{it}{2z}\right|^2\ge\frac12.
\]
For $t>0$ and $\sigma>0$, these points have respectively positive and
negative imaginary parts; the zero endpoints equal $1$.  Thus both line
segments avoid the principal cut.  Along either segment, the derivative
of $(1+\sigma\xi)^{-1/2}$, where $\xi=\pm it/(2z)$, has modulus at most
$2^{3/4}|\xi|/2$.  The mean-value estimate therefore gives
\[
   \left|\left(1\pm\frac{it}{2z}\right)^{-1/2}-1\right|
   \le\frac{2^{3/4}t}{4\rho}.
\]
Using $\int t\,d\mu(t)=1/2$ proves the first estimate in
\eqref{eq:A-uniform-estimates}.  The derivative formula
\eqref{eq:Apm-derivative} and the same geometric bounds give the second.

If $\rho\ge4$, then $\kappa/\rho<1/16$.  The triangle and reverse triangle
inequalities yield $15/16<|A_\pm(z)|<17/16$; in particular, $A_\pm(z)\ne0$.
Finally,
\[
   \left|z\frac{A_\pm'(z)}{A_\pm(z)}\right|
   <\frac1{15}.
\]
\end{proof}

Let
\[
   \mathcal U:=\{z\in\Hplus:A_+(z)A_-(z)\ne0\}.
\]
Since $A_+A_-$ is holomorphic on $\Hplus$, the set $\mathcal U$ is open.
On $\mathcal U$, define
\begin{equation}\label{eq:delta-definition}
   \delta(z):=e^{2i(z-\pi/4)}\frac{A_+(z)}{A_-(z)}.
\end{equation}
Then, throughout $\mathcal U$,
\begin{equation}\label{eq:J0-factorized}
   J_0(z)=\frac{e^{-i(z-\pi/4)}A_-(z)}{\sqrt{2\pi z}}
   \bigl(1+\delta(z)\bigr).
\end{equation}

\begin{proposition}\label{prop:quantitative-separation}
Suppose $J_0(z)=1$ and $z=x+iy=\rho e^{i\theta}$ with
$0<\theta<\pi/4$.  Then $z\in\mathcal U$, and
\begin{equation*}
   y>\frac75,\qquad e^{-2y}<\frac1{16},\qquad
   |\delta(z)|<\frac{17}{240},\qquad
   \rho e^{-2y}<\frac{11}{50}.
\end{equation*}
\end{proposition}

\begin{proof}
By Proposition~\ref{prop:global-localization}, $\rho\ge4$, so
Proposition~\ref{prop:A-uniform-estimates} gives $z\in\mathcal U$.
Using $J_0(z)=1$, \eqref{eq:exact-J0-decomposition}, and
\eqref{eq:A-modulus-bounds}, we obtain
\[
   \sqrt{8\pi}\le\sqrt{2\pi\rho}
   \le\frac{17}{16}(e^y+e^{-y})=\frac{17}{8}\cosh y.
\]
Hence $\cosh y\ge(8/17)\sqrt{8\pi}>11/5$.  Together with
$\cosh(7/5)<11/5$ and $e^{14/5}>16$, this gives $y>7/5$ and
$e^{-2y}<1/16$, since $y>0$.

By \eqref{eq:delta-definition} and \eqref{eq:A-modulus-bounds},
\[
   |\delta(z)|
   =e^{-2y}\left|\frac{A_+(z)}{A_-(z)}\right|
   <\frac1{16}\frac{17}{15}=\frac{17}{240}.
\]
Multiplying the preceding bound for $\sqrt{2\pi\rho}$ by $e^{-y}$ gives
\[
   \sqrt{2\pi\rho}\,e^{-y}
   \le\frac{17}{16}(1+e^{-2y})
   \le\left(\frac{17}{16}\right)^2.
\]
Squaring and using $\pi>3$, we conclude that
\[
   \rho e^{-2y}
   \le\frac{(17/16)^4}{2\pi}
   <\frac{(17/16)^4}{6}
   <\frac{11}{50}.\qedhere
\]
\end{proof}

Proposition~\ref{prop:quantitative-separation} is the point at which the
increasing component becomes uniformly dominant.  The next section shows
that this dominance persists along the whole outward ray.

\section{Radial escape}\label{sec:escape}

\begin{proposition}
\label{prop:principal-radial-escape}
Suppose $J_0(z_1)=1$ and
$z_1=x_1+iy_1=\rho_1e^{i\theta}$ with $\rho_1>0$ and
$0<\theta<\pi/4$.  For
$\tau\ge1$, put $z_\tau=\tau z_1$.  Then
\begin{equation}\label{eq:radial-log-derivative}
   \tau\frac{d}{d\tau}\log|J_0(z_\tau)|>\frac7{25}
   \qquad(\tau\ge1),
\end{equation}
and hence
\begin{equation}\label{eq:principal-radial-escape}
   |J_0(cz_1)|>c^{7/25}\qquad(c>1).
\end{equation}
\end{proposition}

\begin{proof}
Let $\rho_\tau=|z_\tau|=\tau\rho_1$ and
$y_\tau=\Im z_\tau=\tau y_1$.  The argument is fixed and
$\rho_\tau\ge\rho_1\ge4$, so Proposition~\ref{prop:A-uniform-estimates}
applies along the entire ray.

By Proposition~\ref{prop:quantitative-separation},
$\rho_1e^{-2y_1}<11/50$.  Since
\[
   \rho_\tau e^{-2y_\tau}
   =\rho_1e^{-2y_1}\,\tau e^{-2(\tau-1)y_1}
\]
and $\log\tau\le\tau-1$, we have
$\tau e^{-2(\tau-1)y_1}\le1$.  Hence
\begin{equation*}
   \rho_\tau e^{-2y_\tau}<\frac{11}{50}.
\end{equation*}
Also $e^{-2y_\tau}<1/16$.  Therefore
\begin{equation}\label{eq:delta-propagated}
   |\delta(z_\tau)|<\frac{17}{240},
   \qquad
   \rho_\tau|\delta(z_\tau)|<\frac{187}{750}.
\end{equation}
Thus $1+\delta(z_\tau)\ne0$, while
$A_-(z_\tau)\ne0$ by \eqref{eq:A-modulus-bounds}.  The factorization
\eqref{eq:J0-factorized} implies
\begin{equation}\label{eq:J0-ray-nonzero}
   J_0(z_\tau)\ne0\qquad(\tau\ge1).
\end{equation}

Since $J_0$ is entire, \eqref{eq:J0-ray-nonzero} and the identity
$\log|g|=\tfrac12\log|g|^2$ for $g\ne0$ show that $\log|J_0(\tau z_1)|$ is
differentiable near each $\tau\ge1$, including $\tau=1$.
Logarithmic differentiation of \eqref{eq:J0-factorized} along the ray gives
\begin{equation}\label{eq:radial-derivative-decomposition}
\begin{aligned}
   \tau\frac{d}{d\tau}\log|J_0(z_\tau)|
   =y_\tau-\frac12 +\Re\left(
   z_\tau\frac{A_-'(z_\tau)}{A_-(z_\tau)}
   +z_\tau\frac{\delta'(z_\tau)}{1+\delta(z_\tau)}
   \right).
\end{aligned}
\end{equation}
Differentiation of \eqref{eq:delta-definition} gives
\begin{equation}\label{eq:delta-derivative}
   z\delta'(z)=\delta(z)\left(
   2iz+z\frac{A_+'(z)}{A_+(z)}
   -z\frac{A_-'(z)}{A_-(z)}
   \right).
\end{equation}
Using \eqref{eq:delta-derivative}, \eqref{eq:A-modulus-bounds},
\eqref{eq:delta-propagated}, and the triangle inequality, we obtain
\[
   \left|z_\tau\frac{A_-'(z_\tau)}{A_-(z_\tau)}
   +z_\tau\frac{\delta'(z_\tau)}{1+\delta(z_\tau)}\right|
   <\frac1{15}
   +\frac{2(187/750)+(2/15)(17/240)}{1-17/240}
   <\frac{31}{50}.
\]
Substitution into \eqref{eq:radial-derivative-decomposition}, together with
$y_\tau\ge y_1>7/5$, proves \eqref{eq:radial-log-derivative}.
Finally, put
\[
   H(\tau):=\log|J_0(\tau z_1)|-\frac7{25}\log\tau.
\]
The bound \eqref{eq:radial-log-derivative} gives $H'(\tau)>0$ for
$\tau\ge1$.  Since $H(1)=0$, the mean value theorem yields $H(c)>0$ for
$c>1$, which proves \eqref{eq:principal-radial-escape}.
\end{proof}

\begin{theorem}\label{thm:global-radial-escape}
If $z_0\in\C\setminus\{0\}$ and $J_0(z_0)=1$, then
\begin{equation}\label{eq:global-radial-escape}
   |J_0(cz_0)|>c^{7/25}\qquad(c>1).
\end{equation}
\end{theorem}

\begin{proof}
The defining series gives $J_0(-z)=J_0(z)$ and
$J_0(\overline z)=\overline{J_0(z)}$.  Write $z_0=x+iy$.  By
Proposition~\ref{prop:global-localization}, $0<|y|<|x|$.  Thus
$w=|x|+i|y|$ lies in $0<\arg w<\pi/4$ and is one of
$\pm z_0,\pm\overline{z_0}$.  Hence $J_0(w)=1$ and, for every real $c>0$,
$|J_0(cw)|=|J_0(cz_0)|$.  Proposition~\ref{prop:principal-radial-escape}
applied to $w$ proves \eqref{eq:global-radial-escape}.
\end{proof}

\begin{corollary}\label{cor:ray-uniqueness}
For every $\theta\in[0,2\pi)$,
$
   \#\{r>0:J_0(re^{i\theta})=1\}\le1.
$
\end{corollary}

\begin{proof}
If $0<r_1<r_2$ gave two solutions on the same ray, then with
$c=r_2/r_1>1$ and $z_1=r_1e^{i\theta}$,
Theorem~\ref{thm:global-radial-escape} would give
$|J_0(r_2e^{i\theta})|=|J_0(cz_1)|>1$, a contradiction.
\end{proof}

\section{Two-radius consequences}\label{sec:two-radius}

For $R>0$, let $B_R=\{x\in\R^2:|x|<R\}$.  If
$f\in L^1_{\mathrm{loc}}(B_R)$ and $0<r<R$, define the normalized circular mean
\[
   M_rf(x):=\frac1{2\pi}\int_0^{2\pi}
   f\bigl(x+r(\cos t,\sin t)\bigr)\,dt
\]
for almost every $x\in B_{R-r}$.  Fubini's theorem shows that this defines
an $L^1_{\mathrm{loc}}(B_{R-r})$ function independently of the chosen
representative of $f$.

The full-space conditional two-radius theorem goes back to Delsarte and
Lions \cite{DelsarteLions1959,Delsarte1961}.  Volchkov later obtained a
sharp local formulation that separates the spectral obstruction from the
geometric condition on $R,r_1,r_2$ \cite[Theorem~1]{Volchkov1996}.  In
dimension two, define
\[
   Z:=\{z\in\C\setminus\{0\}:J_0(z)=1\},
   \qquad
   \mathcal A:=\left\{\frac{\alpha}{\beta}:\alpha,\beta\in Z\right\}.
\]

\begin{proposition}\label{prop:spectral-equivalence}
For positive radii $r_1,r_2$, the following are equivalent:
\begin{enumerate}
\item $r_1/r_2\in\mathcal A$;
\item there exists $\lambda\in\C\setminus\{0\}$ such that
\[
   J_0(i\lambda r_1)=J_0(i\lambda r_2)=1.
\]
\end{enumerate}
\end{proposition}

\begin{proof}
If $\alpha/\beta=r_1/r_2$ with $\alpha,\beta\in Z$, set
$\lambda=-i\alpha/r_1=-i\beta/r_2$.  Conversely, a common nonzero spectral
point gives $\alpha=i\lambda r_1$, $\beta=i\lambda r_2$, and
$\alpha/\beta=r_1/r_2$.
\end{proof}

A common nonzero spectral parameter gives the following counterexample.

\begin{proposition}\label{prop:explicit-counterexample}
If the equivalent conditions in Proposition~\ref{prop:spectral-equivalence}
hold, then there is a real-valued $C^\infty$ function on $\R^2$ that is not
harmonic and satisfies both circular mean identities.
\end{proposition}

\begin{proof}
For $F_\lambda(x_1,x_2)=e^{\lambda x_1}$,
\[
   M_rF_\lambda
   =F_\lambda\frac1{2\pi}\int_0^{2\pi}e^{\lambda r\cos t}\,dt
   =J_0(i\lambda r)F_\lambda.
\]
Thus the two common spectral identities imply
$M_{r_1}F_\lambda=M_{r_2}F_\lambda=F_\lambda$, whereas
$\Delta F_\lambda=\lambda^2F_\lambda$.  At the origin this Laplacian equals
$\lambda^2\ne0$, so at least one of $\Re F_\lambda$ and $\Im F_\lambda$
is a real-valued smooth nonharmonic example satisfying the same identities.
\end{proof}

Theorem~\ref{thm:global-radial-escape} gives
\begin{equation}\label{eq:positive-ratio-removal}
   (\mathcal A\cap(0,\infty))\setminus\{1\}=\varnothing.
\end{equation}
Indeed, if $\alpha/\beta=c>0$ with $c\ne1$, then either $c>1$ or
$c^{-1}>1$, and two distinct points of $Z$ would lie on one open ray.

We use the following dimension-two specialization of Volchkov's local
theorem \cite[Theorem~1(i)]{Volchkov1996}.

\begin{theorem}\label{thm:Volchkov-local}
Let $r_1,r_2>0$, let $R>r_1+r_2$, and let
$f\in L^1_{\mathrm{loc}}(B_R)$.  Suppose
\[
   M_{r_j}f=f\quad\text{a.e. on }B_{R-r_j},
   \qquad j=1,2.
\]
If $r_1/r_2\notin\mathcal A$, then $\Delta f=0$ in $\D'(B_R)$.
\end{theorem}

\begin{corollary}
\label{cor:local-two-radius}
Let $r_1\ne r_2$ be positive, let $R>r_1+r_2$, and let
$f\in L^1_{\mathrm{loc}}(B_R)$.  If
\[
   M_{r_j}f=f\quad\text{a.e. on }B_{R-r_j},
   \qquad j=1,2,
\]
then $\Delta f=0$ in $\D'(B_R)$.
\end{corollary}

\begin{proof}
The ratio $r_1/r_2$ is positive and different from $1$.  By
\eqref{eq:positive-ratio-removal}, it does not belong to $\mathcal A$.
The conclusion follows from Theorem~\ref{thm:Volchkov-local}.
\end{proof}

\begin{corollary}
\label{cor:global-two-radius}
Let $f\in L^1_{\mathrm{loc}}(\R^2)$, and let $r_1\ne r_2$ be positive.  If
\[
   M_{r_1}f=M_{r_2}f=f\quad\text{a.e. on }\R^2,
\]
then $\Delta f=0$ in $\D'(\R^2)$.  Consequently, $f$ agrees almost
everywhere with a smooth harmonic function.  If $f$ is continuous, then
$f$ itself is harmonic and has the circular mean-value property at every
positive radius.
\end{corollary}

\begin{proof}
For every $R>r_1+r_2$, Corollary~\ref{cor:local-two-radius} applies to the
restriction of $f$ to $B_R$.  Hence $\Delta f=0$ in every such ball and
therefore in $\D'(\R^2)$.  Weyl's lemma gives a smooth harmonic
representative.  The final assertion follows from continuity and the usual
mean-value theorem for harmonic functions.
\end{proof}

\section*{Declarations}

\noindent\textbf{Use of generative AI.}
The Lean formalization was completed using the local Prove2Me
workflow~\cite{Prove2Me} with GPT-6.1 Sol. Generative AI tools also
assisted with mathematical exploration, language, and \LaTeX{} editing.
The authors reviewed the resulting material and take responsibility for
the mathematical content and its correspondence with the formal statements.

\begin{thebibliography}{99}
\raggedright

\bibitem{Delsarte1961}
J.~Delsarte,
\emph{Lectures on Topics in Mean Periodic Functions and the Two-Radius Theorem},
Tata Institute of Fundamental Research Lectures on Mathematics, No.~22,
Tata Institute of Fundamental Research, Bombay, 1961.
\url{https://mathweb.tifr.res.in/Documents/Publications/Lectures/tifr22.pdf}.

\bibitem{DelsarteLions1959}
J.~Delsarte and J.-L.~Lions,
Moyennes g\'en\'eralis\'ees,
\emph{Comment. Math. Helv.} \textbf{33} (1959), 59--69.
\href{https://doi.org/10.1007/BF02565907}{doi:10.1007/BF02565907}.

\bibitem{DLMF}
F.~W.~J. Olver, A.~B. Olde Daalhuis, D.~W. Lozier, B.~I. Schneider,
R.~F. Boisvert, C.~W. Clark, B.~R. Miller, B.~V. Saunders, H.~S. Cohl,
and M.~A. McClain (eds.),
\emph{NIST Digital Library of Mathematical Functions}, Release 1.2.7,
15 June 2026.
\url{https://dlmf.nist.gov/}.

\bibitem{HaymanLingham2019}
W.~K.~Hayman and E.~F.~Lingham,
\emph{Research Problems in Function Theory: Fiftieth Anniversary Edition},
Problem Books in Mathematics, Springer, Cham, 2019.
\href{https://doi.org/10.1007/978-3-030-25165-9}{doi:10.1007/978-3-030-25165-9}.

\bibitem{Volchkov1996}
V.~V.~Volchkov,
The final version of the mean value theorem for harmonic functions,
\emph{Math. Notes} \textbf{59} (1996), no.~3, 247--252.
\href{https://doi.org/10.1007/BF02308535}{doi:10.1007/BF02308535}.

\bibitem{Lean4}
L.~de Moura and S.~Ullrich,
The Lean 4 theorem prover and programming language,
in \emph{Automated Deduction---CADE 28},
Lecture Notes in Comput. Sci., Vol.~12699, Springer, 2021, pp.~625--635.
\href{https://doi.org/10.1007/978-3-030-79876-5_37}{doi:10.1007/978-3-030-79876-5\_37}.

\bibitem{mathlib}
The mathlib Community,
The Lean mathematical library,
in \emph{Proceedings of the 9th ACM SIGPLAN International Conference on
Certified Programs and Proofs}, ACM, 2020, pp.~367--381.
\href{https://doi.org/10.1145/3372885.3373824}{doi:10.1145/3372885.3373824}.

\bibitem{Prove2Me}
S.~Chen, K.~Marwaha, X.~Lu, H.~Yuen, and T.~Peng,
\emph{Prove2Me: An Open Collaborative Platform for Scaling Math Formalization},
arXiv:2608.28433, 2026.
\url{https://doi.org/10.48550/arXiv.2608.28433}.

\bibitem{Ferudun2026}
A.~Ferudun,
\emph{The Equation $J_0(z)=1$ Has at Most One Solution on Each Ray:
An Answer to a Question of Zalcman},
EulerSolve Research Papers, AMR-022-2045, preprint,
3 October 2026.
\url{https://eulersolve.org/papers/amr-022-2045/}.

\end{thebibliography}

\end{document}
